%% kapa: exam=1 round=31 year=2020 period=1 form=A subject=civil_law subject_ko=민법 seq=1 number=1
%% source: Q-Net 감정평가사 기출문제 (cst00309) artlSeq=5207219 file=extracted/1008/5207219/A형/A형/2020년도 제31회 감정평가사 1차 시험_1교시_A형.pdf
%% conversion: kordoc 4.24.51 -> tools/gen1.py (auto); review=proofread-sonnet; flags=none
\begin{question}{1}
법원(法源)에 관한 설명으로 옳지 않은 것은? (다툼이 있으면 판례에 따름)
\begin{choices}
  \item 사회구성원이 관습법으로 승인된 관행의 법적 구속력을 확신하지 않게 된 때에는 그 관습법은 효력을 잃는다.
  \item 헌법의 기본권은 특별한 사정이 없으면 사법관계에 직접 적용된다.
  \item 법원은 당사자의 주장․증명을 기다림이 없이 관습법을 직권으로 조사ㆍ확정하여야 한다.
  \item 우리나라가 가입한 국제조약은 일반적으로 민법이나 상법 또는 국제사법보다 우선적으로 적용된다.
  \item 관습법은 법령에 저촉되지 아니하는 한 법칙으로서의 효력이 있다.
\end{choices}
\end{question}
\answer{2}
